本研究針對澳門氹仔地區的共享單車系統，提出了一個綜合優化模型，旨在解決站點選址與初始投放量優化以及再平衡問題。通過將氹仔劃分為居民區、旅遊區和樞紐區，並考慮空間異質性、時間不對稱性和資源有限性等因素，模型採用混合整數規劃方法，以最小化服務缺口為目標。在再平衡方面，基於有向圖的調度模型整合了運輸成本與懲罰成本，支持多輛卡車參與調度。數值試驗驗證了模型的有效性，為共享單車系統的規劃與運營提供了科學依據。

\begin{center}
    \textbf{關鍵字: 共享單車, 調度模型}
\end{center}

